% =============================================================================
%  1.3 Async/Await
% =============================================================================
\section{Async/Await}

\term{Асинхронные функции}{Async Functions} — синтаксический сахар над
промисами, позволяющий писать асинхронный код в синхронном стиле.
Ключевое слово \inl{async} превращает функцию в генератор промисов,
а \inl{await} приостанавливает выполнение до разрешения промиса.

Под капотом среда исполнения использует конечный автомат: каждый
\inl{await} создаёт точку приостановки, и функция возобновляется
через очередь микрозадач.

\subsection{Базовый синтаксис}

\begin{code}{javascript}
async function loadDashboard(userId) {
  try {
    const user = await fetchUser(userId);
    const posts = await fetchPosts(user.id);
    return { user, posts };
  } catch (err) {
    console.error('Ошибка загрузки:', err);
    throw err;
  }
}
\end{code}

\subsection{Параллельное выполнение}

\subsubsection{Антипаттерн: последовательный await}

Время выполнения последовательных \inl{await} складывается линейно:
$T = \sum_{i=1}^{n} t_i$, тогда как при параллельном запуске
$T = \max(t_1, \ldots, t_n)$.

\begin{equation}
  \text{Speedup} = \frac{\sum_{i=1}^{n} t_i}{\max_{i} t_i}
\end{equation}

\begin{code}{javascript}
// Плохо — последовательно
const user = await fetchUser(1);
const posts = await fetchPosts(1);   // ждёт завершения fetchUser

// Хорошо — параллельно
const [user, posts] = await Promise.all([
  fetchUser(1),
  fetchPosts(1),
]);
\end{code}

\note{Используйте \inl{Promise.all} внутри \inl{async}-функций,
когда операции не зависят друг от друга.  Это особенно важно при
множественных запросах к \abbr{API}{Application Programming Interface}{программный интерфейс}.}

\important{Не забывайте обрабатывать ошибки в \inl{async}-функциях.
Необработанный \inl{reject} в Node.js приводит к завершению процесса
начиная с версии~15.}
